% ============================================================
%  E-Descarte: Plataforma Digital de Gestão e Denúncia de
%  Resíduos Eletroeletrônicos com Módulo Legal Integrado
%  Artigo científico – LaTeX / ABNT – Overleaf-ready
% ============================================================

\documentclass[12pt, a4paper]{article}

% ── Codificação e Idioma ─────────────────────────────────────
\usepackage[utf8]{inputenc}
\usepackage[T1]{fontenc}
\usepackage[brazil]{babel}

% ── Geometria da Página (ABNT NBR 14724) ────────────────────
\usepackage[
  top=3cm, bottom=2cm,
  left=3cm, right=2cm
]{geometry}

% ── Tipografia ───────────────────────────────────────────────
\usepackage{times}            % Times New Roman
\usepackage{microtype}        % Melhorias tipográficas
\usepackage{setspace}         % Controle de espaçamento
\usepackage{indentfirst}      % Recuo no primeiro parágrafo

% ── Cabeçalho e Rodapé ───────────────────────────────────────
\usepackage{fancyhdr}
\pagestyle{fancy}
\fancyhf{}
\fancyhead[R]{\small\thepage}
\fancyhead[L]{\small\textit{E-Descarte: Plataforma Digital de Gestão de Resíduos Eletroeletrônicos}}
\renewcommand{\headrulewidth}{0.4pt}

% ── Títulos de Seção ─────────────────────────────────────────
\usepackage{titlesec}
\titleformat{\section}{\normalfont\normalsize\bfseries\uppercase}{}{0em}{}
\titleformat{\subsection}{\normalfont\normalsize\bfseries}{}{0em}{}
\titleformat{\subsubsection}{\normalfont\normalsize\itshape}{}{0em}{}

% ── Figuras, Tabelas e Listagens ────────────────────────────
\usepackage{graphicx}
\usepackage{float}
\usepackage{caption}
\usepackage{subcaption}
\usepackage{booktabs}
\usepackage{array}
\usepackage{tabularx}
\usepackage{longtable}
\usepackage{multirow}
\usepackage{listings}
\usepackage{xcolor}

% Estilo de código-fonte
\lstset{
  basicstyle=\footnotesize\ttfamily,
  breaklines=true,
  frame=single,
  rulecolor=\color{gray!50},
  backgroundcolor=\color{gray!8},
  numbers=left,
  numberstyle=\tiny\color{gray},
  keywordstyle=\color{blue!70},
  commentstyle=\color{green!50!black},
  stringstyle=\color{orange!80!black},
  captionpos=b,
  tabsize=2,
  showstringspaces=false
}

% ── Diagramas TikZ ───────────────────────────────────────────
\usepackage{tikz}
\usetikzlibrary{shapes.geometric, arrows.meta, positioning, fit, backgrounds}

% ── Matemática ───────────────────────────────────────────────
\usepackage{amsmath}
\usepackage{amssymb}

% ── URLs e Hiperlinks ────────────────────────────────────────
\usepackage{url}
\usepackage[hidelinks]{hyperref}

% ── Referências ABNT ────────────────────────────────────────
\usepackage[alf, abnt-etal-list=0]{abntex2cite}

% ── Espaçamento entre parágrafos e linhas ───────────────────
\setlength{\parindent}{1.25cm}
\setlength{\parskip}{0pt}
\onehalfspacing

% ── Numeração de equações ────────────────────────────────────
\numberwithin{equation}{section}
\numberwithin{figure}{section}
\numberwithin{table}{section}

% ============================================================
%  INÍCIO DO DOCUMENTO
% ============================================================
\begin{document}

% ── Página de Título ─────────────────────────────────────────
\begin{titlepage}
  \begin{center}
    {\large \textbf{UNIVERSIDADE JORGE AMADO -- UNIJORGE}}\\[0.3cm]
    {\large \textbf{CURSO DE CIÊNCIA DA COMPUTAÇÃO}}\\[3cm]

    {\Large \textbf{E-DESCARTE: PLATAFORMA DIGITAL DE GESTÃO E DENÚNCIA DE
    RESÍDUOS ELETROELETRÔNICOS COM MÓDULO LEGAL INTEGRADO}}\\[3cm]

    {\normalsize
      Renoir Auguste Santos Silva e Colaboradores\\[0.3cm]
      Disciplina: EcoTech -- Observatório Digital de Direito Ambiental\\[3cm]
    }

    \vfill
    {\normalsize Salvador (BA)\\2025}
  \end{center}
\end{titlepage}

\clearpage
\setcounter{page}{2}

% ── Resumo ───────────────────────────────────────────────────
\begin{center}
  {\normalsize\bfseries
  E-DESCARTE: PLATAFORMA DIGITAL DE GESTÃO E DENÚNCIA DE RESÍDUOS
  ELETROELETRÔNICOS COM MÓDULO LEGAL INTEGRADO\\[0.5cm]}
\end{center}

\begin{center}
  {\normalsize\bfseries RESUMO}
\end{center}
\noindent
O crescimento acelerado da indústria tecnológica tem como consequência direta
o aumento exponencial na geração de resíduos eletroeletrônicos (REEE), tornando
essa problemática um dos maiores desafios ambientais contemporâneos. O Brasil,
quarto maior gerador mundial de lixo eletrônico, produz aproximadamente 2,1
milhões de toneladas anuais, volume que, descartado de forma inadequada, libera
metais pesados como chumbo, mercúrio e cádmio no solo, na água e no ar.
Paralelamente, a legislação brasileira conta com instrumentos robustos para
o enfrentamento desse problema, especialmente a Lei nº 12.305/2010 (Política
Nacional de Resíduos Sólidos -- PNRS), o Decreto nº 10.240/2020 (logística
reversa de eletroeletrônicos) e a Resolução CONAMA 401/2008, cujo cumprimento
permanece insatisfatório na maioria dos municípios brasileiros. Este artigo
apresenta o \textit{E-Descarte}, uma plataforma digital desenvolvida como
resposta tecnológica a esse cenário. O sistema integra três módulos funcionais:
(i) um mapa interativo de pontos de coleta licenciados com geolocalização via
API do Google Maps; (ii) um sistema de denúncias de descartes irregulares com
suporte a anonimato, georreferenciamento e envio de fotografias; e (iii) uma
calculadora de impacto ambiental estimado. Como diferencial técnico e pedagógico,
a plataforma conta com um Módulo Legal Integrado que apresenta, em tempo real,
a legislação ambiental aplicável ao tipo de resíduo informado pelo usuário,
contribuindo para a conscientização cidadã e para o controle social da aplicação
das leis ambientais. A arquitetura do sistema adota uma abordagem de separação
de responsabilidades em três camadas: \textit{frontend} React (Web) e React
Native (móvel), \textit{backend} Node.js/Express com autenticação JWT, e
persistência híbrida em PostgreSQL e MongoDB. Os resultados esperados incluem
o aumento dos índices de descarte correto, a ampliação do controle social sobre
descartes irregulares e o fortalecimento da consciência legal ambiental da
população.

\vspace{0.5cm}
\noindent
\textbf{Palavras-chave}: Resíduos Eletroeletrônicos. Lixo Eletrônico. Política
Nacional de Resíduos Sólidos. Plataforma Digital Ambiental. Geolocalização.
Node.js. React. Módulo Legal Integrado.

\clearpage

% ── Abstract ─────────────────────────────────────────────────
\begin{center}
  {\normalsize\bfseries ABSTRACT}
\end{center}
\noindent
The accelerated growth of the technology industry has as a direct consequence
the exponential increase in the generation of waste electrical and electronic
equipment (WEEE), making this issue one of the greatest contemporary
environmental challenges. Brazil, the world's fourth largest generator of
electronic waste, produces approximately 2.1 million tons per year, a volume
that, when improperly disposed of, releases heavy metals such as lead, mercury
and cadmium into the soil, water and air. At the same time, Brazilian legislation
has robust instruments to address this problem, especially Law No.\ 12,305/2010
(National Solid Waste Policy -- PNRS), Decree No.\ 10,240/2020 (reverse
logistics of electronic equipment) and CONAMA Resolution 401/2008, whose
compliance remains unsatisfactory in most Brazilian municipalities. This article
presents \textit{E-Descarte}, a digital platform developed as a technological
response to this scenario. The system integrates three functional modules: (i) an
interactive map of licensed collection points with geolocation via the Google
Maps API; (ii) an irregular disposal reporting system with support for anonymity,
georeferencing and photo submission; and (iii) an estimated environmental impact
calculator. As a technical and pedagogical differentiator, the platform features
an Integrated Legal Module that displays, in real time, the environmental
legislation applicable to the type of waste reported by the user, contributing
to citizen awareness and social control of the enforcement of environmental laws.
The system architecture adopts a three-layer separation of concerns approach:
React front-end (Web) and React Native (mobile), Node.js/Express back-end with
JWT authentication, and hybrid persistence in PostgreSQL and MongoDB.

\vspace{0.5cm}
\noindent
\textbf{Keywords}: Electronic Waste. WEEE. National Solid Waste Policy.
Digital Environmental Platform. Geolocation. Node.js. React.
Integrated Legal Module.

\clearpage

% ── Sumário ──────────────────────────────────────────────────
\tableofcontents
\clearpage

% ============================================================
\section{INTRODUÇÃO}
% ============================================================

A digitalização acelerada da sociedade contemporânea trouxe consigo benefícios
inegáveis para a produtividade, a comunicação e o acesso à informação. Contudo,
esse processo também gerou um problema ambiental silencioso e crescente:
a proliferação de resíduos eletroeletrônicos (REEE), popularmente denominados
``lixo eletrônico''. Celulares obsoletos, computadores descartados, baterias
desgastadas e televisores substituídos acumulam-se em lixões, terrenos baldios
e aterros convencionais em ritmo que supera em muito a capacidade de gestão dos
sistemas públicos de saneamento \cite{abrelpe2023}.

Conforme dados do Relatório Global sobre Resíduos Eletrônicos publicado pela
ONU, o mundo gerou aproximadamente 53,6 milhões de toneladas métricas de
resíduos eletrônicos em 2019, com crescimento projetado para 74,7 milhões de
toneladas até 2030 \cite{forti2020}. O Brasil ocupa a quarta posição nesse
ranking, produzindo cerca de 2,1 milhões de toneladas por ano e reciclando
formalmente menos de 3\% desse volume \cite{abrelpe2023}.

A gravidade do problema não reside apenas no volume gerado, mas na composição
química dos equipamentos descartados. Dispositivos eletrônicos contêm substâncias
altamente tóxicas como chumbo, cádmio, mercúrio, arsênio, cromo hexavalente e
retardantes de chama bromados. Quando lançados em aterros convencionais ou
queimados a céu aberto, esses compostos percolam para os lençóis freáticos,
contaminam o solo por décadas e emitem gases e partículas nocivas à saúde
humana, causando danos neurológicos, renais e carcinogênicos
\cite{baldé2017, singh2020}.

O ordenamento jurídico brasileiro oferece resposta normativa expressiva a esse
desafio. A Lei nº 12.305/2010, que institui a Política Nacional de Resíduos
Sólidos (PNRS), estabelece o princípio da responsabilidade compartilhada pelo
ciclo de vida dos produtos e cria o instrumento da logística reversa,
obrigando fabricantes, importadores e distribuidores a recolher e dar destinação
adequada aos resíduos de seus produtos \cite{brasil2010pnrs}. O Decreto nº
10.240/2020 regulamentou especificamente a logística reversa de produtos
eletroeletrônicos de uso doméstico, definindo metas e prazos para a estruturação
de pontos de coleta em todo o território nacional \cite{brasil2020decreto}. A
Resolução CONAMA 401/2008, por sua vez, disciplina o gerenciamento de pilhas
e baterias \cite{conama401}. Apesar desse arcabouço legal, a fiscalização é
fragmentada e o desconhecimento da legislação por parte da população e das
próprias empresas é generalizado \cite{moreira2019}.

Nesse contexto, este trabalho apresenta o \textit{E-Descarte}, uma plataforma
digital desenvolvida no âmbito da disciplina EcoTech -- Observatório Digital de
Direito Ambiental, do Curso de Ciência da Computação da Universidade Jorge Amado
(UNIJORGE). O sistema propõe uma solução tecnológica integrada que une
geolocalização de pontos de coleta, registro georreferenciado de denúncias e,
de forma inovadora, um Módulo Legal Integrado que apresenta ao cidadão, em
tempo real, a legislação ambiental aplicável ao tipo de resíduo que deseja
reportar.

\subsection{Objetivo Geral}

Desenvolver uma plataforma digital completa -- composta por aplicação Web,
aplicativo móvel e API REST -- para a gestão e denúncia de resíduos
eletroeletrônicos no contexto brasileiro, integrando funcionalidades de
geolocalização, controle social e educação legal ambiental.

\subsection{Objetivos Específicos}

\begin{enumerate}
  \item Mapear e disponibilizar, em interface cartográfica interativa, os pontos
        de coleta de resíduos eletroeletrônicos licenciados;
  \item Implementar um sistema de denúncias de descartes irregulares com suporte
        a anonimato, georreferenciamento GPS e envio de evidências fotográficas;
  \item Desenvolver uma calculadora de impacto ambiental que sensibilize o
        usuário sobre os benefícios do descarte correto;
  \item Criar um Módulo Legal Integrado que apresente, dinamicamente, a
        legislação ambiental brasileira aplicável conforme o tipo de resíduo
        informado;
  \item Construir uma arquitetura de \textit{software} escalável, segura e bem
        documentada, seguindo boas práticas de desenvolvimento \textit{full-stack}.
\end{enumerate}

\subsection{Justificativa}

A relevância deste projeto fundamenta-se em três eixos complementares.
Do ponto de vista ambiental, o descarte inadequado de eletrônicos representa
uma das formas mais graves de contaminação por metais pesados no Brasil,
com impactos diretos sobre a saúde pública e os ecossistemas. Do ponto de
vista legal, existe um hiato significativo entre a robustez da legislação
ambiental brasileira e a efetividade de sua aplicação, lacuna que a tecnologia
pode ajudar a reduzir por meio da informação e do controle social. Do ponto
de vista tecnológico, a combinação de aplicações móveis, APIs geoespaciais e
bases de dados híbridas oferece um conjunto de ferramentas adequado e atual
para o enfrentamento desse problema de forma escalável e acessível.

% ============================================================
\section{FUNDAMENTAÇÃO TEÓRICA}
% ============================================================

\subsection{Resíduos Eletroeletrônicos: Conceito e Dimensão do Problema}

Os resíduos de equipamentos elétricos e eletrônicos (REEE) abrangem qualquer
produto que funcione com eletricidade ou campo eletromagnético e que tenha
chegado ao fim de sua vida útil. A classificação da Diretiva Europeia
2012/19/EU, amplamente adotada na literatura científica, organiza os REEE em
dez categorias, incluindo grandes e pequenos eletrodomésticos, equipamentos de
tecnologia da informação e comunicação (TIC), equipamentos de consumo e
iluminação, ferramentas elétricas, brinquedos eletrônicos e equipamentos
médicos \cite{europarlamento2012}.

Do ponto de vista da composição material, os equipamentos eletrônicos são
montagens complexas que integram plásticos (até 30\% da massa), metais ferrosos
e não ferrosos (ferro, cobre, alumínio -- até 60\% da massa) e substâncias
perigosas em concentrações variadas \cite{baldé2017}. O chumbo, presente em
soldas e baterias chumbo-ácido, é um neurotóxico que compromete o desenvolvimento
cognitivo de crianças mesmo em baixas concentrações. O mercúrio, encontrado em
displays de cristal líquido e baterias de botão, é um potente nefrotóxico
e neurotóxico. O cádmio, presente em baterias recarregáveis e em componentes
de soldagem, é classificado como carcinogênico para humanos pela Agência
Internacional de Pesquisa sobre Câncer (IARC) \cite{singh2020}.

\subsection{Cenário Brasileiro de Geração e Gestão de REEE}

O Brasil apresenta um paradoxo característico dos países emergentes: é um dos
maiores mercados consumidores de eletrônicos do mundo -- fator impulsionado pelo
crescimento da classe média nas décadas de 2000 e 2010 -- mas dispõe de
infraestrutura de coleta e reciclagem formal ainda incipiente \cite{abrelpe2023}.

A Associação Brasileira de Limpeza Pública e Resíduos Especiais (ABRELPE)
estimou, em seu Panorama dos Resíduos Sólidos no Brasil de 2022, que apenas
cerca de 3\% dos resíduos eletrônicos gerados no país recebem destinação
formalmente adequada. O restante é descartado junto ao lixo doméstico comum,
em lixões a céu aberto ou vendido para sucateiros informais que frequentemente
recorrem a práticas de processamento primitivas e altamente poluidoras, como
a queima de cabos para recuperação do cobre \cite{abrelpe2023}.

A distribuição geográfica da geração de REEE no Brasil é desigual. As regiões
Sudeste e Sul concentram a maior parcela da geração (cerca de 70\% do total
nacional), reflexo da maior densidade demográfica e do maior poder aquisitivo
médio. Contudo, a carência de infraestrutura de coleta e reciclagem é mais
grave justamente nas regiões Norte e Nordeste, onde a proporção de municípios
sem qualquer ponto de coleta formal é mais elevada \cite{moreira2019}.

\subsection{Arcabouço Legal Brasileiro para Resíduos Eletroeletrônicos}

O principal instrumento normativo para o enfrentamento dos REEE no Brasil é
a Lei nº 12.305, de 02 de agosto de 2010, que institui a Política Nacional de
Resíduos Sólidos (PNRS). Essa lei estabelece os princípios, objetivos e
instrumentos para a gestão integrada de resíduos sólidos no país, adotando
como eixo central o conceito de responsabilidade compartilhada pelo ciclo de
vida dos produtos -- princípio pelo qual fabricantes, importadores,
distribuidores, comerciantes, consumidores e titulares dos serviços de limpeza
urbana respondem solidariamente pela destinação final dos resíduos que geram
\cite{brasil2010pnrs}.

O instrumento operacional central da PNRS para os REEE é a logística reversa,
definida no Art. 3º, XII, como o ``instrumento de desenvolvimento econômico e
social caracterizado por um conjunto de ações, procedimentos e meios destinados
a viabilizar a coleta e a restituição dos resíduos sólidos ao setor empresarial,
para reaproveitamento, em seu ciclo ou em outros ciclos produtivos''. A obrigação
de estruturar sistemas de logística reversa para produtos eletroeletrônicos
e seus componentes foi regulamentada pelo Decreto nº 10.240, de 12 de fevereiro
de 2020, que estabeleceu metas progressivas de coleta e reciclagem, além de
exigir que fabricantes e importadores tornem públicas as informações sobre os
pontos de coleta \cite{brasil2020decreto}.

A Resolução do Conselho Nacional do Meio Ambiente (CONAMA) nº 401, de 04 de
novembro de 2008, estabelece os limites máximos de chumbo, cádmio e mercúrio
para pilhas e baterias comercializadas no Brasil, bem como os critérios e
padrões para o gerenciamento ambientalmente adequado de seus resíduos,
proibindo o descarte em aterros, terrenos baldios, corpos d'água ou queima
a céu aberto \cite{conama401}.

Complementarmente, a Lei nº 9.605/1998 (Lei de Crimes Ambientais) tipifica
como crime ambiental, nos termos do Art. 54, causar poluição de qualquer
natureza em níveis tais que resultem ou possam resultar em danos à saúde
humana, mortandade de animais ou destruição significativa da flora, prevendo
pena de reclusão de um a quatro anos e multa. O descarte irregular de grandes
volumes de REEE pode enquadrar-se nessa tipificação \cite{brasil1998crimes}.

A Lei Geral de Proteção de Dados Pessoais (LGPD), Lei nº 13.709/2018, introduz
uma dimensão adicional ao problema dos REEE: os dispositivos eletroeletrônicos
frequentemente armazenam dados pessoais sensíveis de seus usuários. A destinação
inadequada de equipamentos sem a prévia eliminação segura desses dados configura
potencial violação à LGPD \cite{brasil2018lgpd}.

\subsection{Tecnologias de Informação Aplicadas à Gestão Ambiental}

A aplicação de tecnologias da informação e comunicação (TIC) à gestão ambiental
constitui campo interdisciplinar em expansão, denominado Green IT ou
\textit{Environmental Informatics} \cite{watson2010}. Sistemas de informações
geográficas (SIG/GIS), plataformas de crowdsourcing ambiental, aplicativos
móveis de monitoramento e plataformas de \textit{big data} para análise de
padrões de descarte são exemplos de soluções que têm sido desenvolvidas e
avaliadas na literatura científica recente.

Aplicativos de mapeamento colaborativo de pontos de coleta, como o iRecycle
(EUA) e o RecycleMap (Europa), demonstraram que a disponibilização de
informações geolocalizadas aumenta significativamente a taxa de descarte
correto entre usuários dessas plataformas \cite{reno2020}. No contexto
brasileiro, iniciativas como o TrashOut e o Coletai apontam na mesma direção,
embora com penetração ainda limitada \cite{ribeiro2021}.

A adoção de arquiteturas de \textit{software} baseadas em microsserviços,
APIs REST e aplicações \textit{mobile-first} tem sido apontada como abordagem
adequada para plataformas de participação cidadã ambiental, dado que permite
escalabilidade, manutenção simplificada e acesso ubíquo via dispositivos móveis
\cite{fowler2014}.

% ============================================================
\section{TRABALHOS RELACIONADOS}
% ============================================================

Diversas iniciativas, tanto no âmbito acadêmico quanto no setor privado,
têm abordado o problema da gestão de REEE por meio de soluções tecnológicas.
Esta seção apresenta uma revisão dos principais trabalhos identificados,
posicionando o E-Descarte em relação a eles.

\subsection{iRecycle e Plataformas Internacionais}

O iRecycle, desenvolvido pela Earth911 nos Estados Unidos, é uma das plataformas
de localização de pontos de coleta de maior penetração no mundo anglófono.
A plataforma disponibiliza um banco de dados com mais de 100.000 locais de
reciclagem e é acessível via \textit{web} e aplicativo móvel. Sua limitação
central reside na ausência de funcionalidades de denúncia e de conteúdo
educativo-legal, sendo estritamente utilitária na função de localização
\cite{earth911}.

O TrashOut, por sua vez, é uma plataforma \textit{mobile} de crowdsourcing
focada no mapeamento de despejos ilegais de resíduos em geral. Disponível em
mais de 180 países, permite que usuários fotografem e geolocalizem pontos de
descarte irregular, criando um banco de dados colaborativo \cite{trashout}.
Em comparação ao E-Descarte, o TrashOut não integra legislação ambiental
específica nem se especializa no contexto regulatório brasileiro.

\subsection{Iniciativas Brasileiras}

\citeonline{ribeiro2021} apresentam uma análise de aplicativos brasileiros
voltados para a gestão de resíduos sólidos, identificando lacunas recorrentes:
ausência de informações sobre a legislação vigente, baixa integração com órgãos
de fiscalização municipais e estaduais, e falta de mecanismos de \textit{feedback}
sobre a resolução das denúncias realizadas. O E-Descarte foi concebido para
suprir especificamente essas lacunas.

\citeonline{moreira2019} propõem um framework conceitual para sistemas de
informação ambiental voltados à gestão de REEE no Brasil, enfatizando a
necessidade de integração entre as dimensões tecnológica, legal e educativa.
O Módulo Legal Integrado do E-Descarte constitui resposta direta a essa
recomendação.

\subsection{Sistemas de Denúncia Ambiental}

O Sistema Nacional de Informações sobre o Meio Ambiente (SINIMA) do IBAMA
e a plataforma GOV.BR mantêm canais de denúncia ambiental, porém com interfaces
pouco amigáveis, ausência de suporte a georreferenciamento automático e
sem integração com sistemas de mapeamento cartográfico \cite{ibama2022}.
A proposta do E-Descarte complementa esses sistemas ao oferecer uma interface
moderna, mobile-first e com feedback imediato ao usuário sobre o embasamento
legal de sua denúncia.

\subsection{Posicionamento do E-Descarte}

A Tabela~\ref{tab:comparativo} apresenta um comparativo entre o E-Descarte e
os principais sistemas relacionados, evidenciando o diferencial do Módulo Legal
Integrado.

\begin{table}[H]
  \centering
  \caption{Comparativo entre o E-Descarte e sistemas relacionados}
  \label{tab:comparativo}
  \begin{tabularx}{\textwidth}{|l|X|X|X|X|X|}
    \hline
    \textbf{Sistema} &
    \textbf{Mapa de coleta} &
    \textbf{Denúncia} &
    \textbf{Calc. Impacto} &
    \textbf{Módulo Legal} &
    \textbf{BR-ABNT} \\
    \hline
    iRecycle       & Sim & Não & Não & Não & Não \\
    \hline
    TrashOut       & Sim & Sim & Não & Não & Não \\
    \hline
    Coletai (BR)   & Sim & Não & Não & Não & Parcial \\
    \hline
    IBAMA/GOV.BR   & Não & Sim & Não & Não & Sim \\
    \hline
    \textbf{E-Descarte} &
    \textbf{Sim} & \textbf{Sim} & \textbf{Sim} & \textbf{Sim} & \textbf{Sim} \\
    \hline
  \end{tabularx}
  \fonte{Elaborado pelos autores (2025).}
\end{table}

% ============================================================
\section{METODOLOGIA}
% ============================================================

O desenvolvimento do E-Descarte seguiu uma abordagem metodológica mista,
combinando elementos da Engenharia de \textit{Software} iterativo-incremental
com os princípios da pesquisa aplicada em Informática Ambiental.

\subsection{Abordagem de Pesquisa e Desenvolvimento}

A pesquisa classifica-se como aplicada, de natureza qualitativa e quantitativa,
com objetivo exploratório-descritivo. Do ponto de vista dos procedimentos
técnicos, constitui-se como pesquisa bibliográfica, documental e de
desenvolvimento de sistema \cite{gil2017}. A revisão bibliográfica abrangeu
publicações científicas indexadas nas bases ACM Digital Library, IEEE Xplore,
SciELO e Google Scholar, com recorte temporal entre 2015 e 2025, utilizando-se
os descritores: ``resíduos eletroeletrônicos'', ``WEEE management system'',
``lixo eletrônico Brasil'', ``environmental mobile application'' e
``logística reversa eletroeletrônicos Brasil''.

A análise da legislação ambiental foi conduzida com base nas fontes primárias
disponíveis no Portal da Legislação do Planalto e nos documentos oficiais do
IBAMA, CONAMA e Ministério do Meio Ambiente.

\subsection{Metodologia de Desenvolvimento de Software}

O desenvolvimento do sistema adotou uma versão simplificada do \textit{framework}
Scrum, adaptada ao contexto acadêmico de grupo pequeno. O processo foi organizado
em quatro \textit{sprints} de duas semanas, conforme descrito na
Tabela~\ref{tab:sprints}.

\begin{table}[H]
  \centering
  \caption{Organização das sprints de desenvolvimento}
  \label{tab:sprints}
  \begin{tabularx}{\textwidth}{|c|X|X|}
    \hline
    \textbf{Sprint} & \textbf{Objetivo} & \textbf{Entregável} \\
    \hline
    1 & Definição de arquitetura, configuração dos ambientes e estrutura
        do banco de dados & Infraestrutura de backend, migrations SQL,
        modelos MongoDB \\
    \hline
    2 & Desenvolvimento da API REST completa (autenticação, CRUD de pontos
        de coleta, denúncias e módulo legal) & API funcional com
        documentação de endpoints \\
    \hline
    3 & Desenvolvimento do frontend web (React/Vite) com todas as páginas
        e integração à API & Aplicação web funcional \\
    \hline
    4 & Desenvolvimento do aplicativo móvel (React Native/Expo),
        testes de integração e ajustes finais & Aplicativo móvel
        funcional e código revisado \\
    \hline
  \end{tabularx}
  \fonte{Elaborado pelos autores (2025).}
\end{table}

\subsection{Modelagem do Sistema}

A modelagem foi conduzida em três níveis: (i) levantamento e especificação de
requisitos funcionais e não funcionais; (ii) modelagem da arquitetura do sistema
por meio de diagrama de camadas e diagrama de componentes; e (iii) modelagem
dos dados por meio do Diagrama Entidade-Relacionamento (DER) para o banco
relacional PostgreSQL e da definição dos esquemas de coleções para o MongoDB.

Os requisitos funcionais foram elicitados a partir da análise de sistemas
relacionados, da revisão da literatura e de sessões de \textit{brainstorming}
da equipe. Os requisitos não funcionais foram definidos com base nos atributos
de qualidade da norma ISO/IEC 25010, com ênfase em usabilidade, segurança
(autenticação e proteção de dados), portabilidade (\textit{web} e móvel) e
manutenibilidade.

% ============================================================
\section{DESENVOLVIMENTO DO PROJETO}
% ============================================================

\subsection{Arquitetura Geral do Sistema}

O E-Descarte adota uma arquitetura em três camadas (\textit{three-tier
architecture}), com clara separação entre a camada de apresentação
(\textit{frontend}), a camada de negócio (\textit{backend}/API) e a camada
de dados (\textit{databases}). Essa escolha arquitetural favorece a
manutenibilidade, a escalabilidade horizontal e a possibilidade de substituição
ou expansão de qualquer camada de forma independente \cite{sommerville2019}.

A Figura~\ref{fig:arquitetura} ilustra a arquitetura geral do sistema.

\begin{figure}[H]
  \centering
  \begin{tikzpicture}[
    node distance=1.2cm and 2cm,
    every node/.style={font=\small},
    box/.style={draw, rounded corners=4pt, minimum width=3.2cm,
                minimum height=0.9cm, align=center, fill=white},
    layer/.style={draw, dashed, rounded corners=6pt, fill=gray!5,
                  inner sep=8pt},
    arr/.style={-{Stealth}, thick, gray}
  ]

  %% Camada Frontend
  \node[box, fill=teal!12] (web)    {Frontend Web\\React/Vite};
  \node[box, fill=teal!12, right=1.6cm of web] (mobile) {App Móvel\\React Native/Expo};

  %% Camada API
  \node[box, fill=blue!10, below=1.5cm of web, xshift=2cm] (api)
        {Backend\\Node.js / Express\\JWT Auth};

  %% Camada de Dados
  \node[box, fill=purple!10, below=1.5cm of api, xshift=-2cm] (pg)
        {PostgreSQL\\(dados estruturados)};
  \node[box, fill=purple!10, below=1.5cm of api, xshift=2cm] (mongo)
        {MongoDB\\(mídia e logs)};

  %% Serviços externos
  \node[box, fill=orange!10, right=1.6cm of api] (gmaps)
        {Google Maps API\\(geocoding)};
  \node[box, fill=orange!10, below=0.7cm of gmaps] (notif)
        {E-mail / SMS\\(Nodemailer/Twilio)};

  %% Setas
  \draw[arr] (web)    -- (api);
  \draw[arr] (mobile) -- (api);
  \draw[arr] (api)    -- (pg);
  \draw[arr] (api)    -- (mongo);
  \draw[arr] (api)    -- (gmaps);
  \draw[arr] (api)    -- (notif);

  \end{tikzpicture}
  \caption{Arquitetura geral do E-Descarte}
  \label{fig:arquitetura}
  \fonte{Elaborado pelos autores (2025).}
\end{figure}

\subsection{Módulo Legal Integrado}

O Módulo Legal Integrado constitui o principal diferencial técnico-pedagógico
do E-Descarte. Ele é implementado no \textit{backend} por meio do módulo
\texttt{legalContent.js}, que mantém uma base de conhecimento estruturada das
legislações ambientais brasileiras aplicáveis a cada categoria de resíduo
eletroeletrônico. Quando o usuário seleciona o tipo de resíduo na tela de
denúncia, o \textit{frontend} realiza uma requisição ao endpoint
\texttt{GET /api/legislacao/tipo/:tipoResiduo}, que retorna o conjunto de
legislações aplicáveis em tempo real.

O Quadro~\ref{quad:legislacao} apresenta o mapeamento entre tipos de resíduo
e a legislação correspondente implementada no módulo.

\begin{table}[H]
  \centering
  \caption{Mapeamento resíduo--legislação no Módulo Legal Integrado}
  \label{quad:legislacao}
  \begin{tabularx}{\textwidth}{|l|X|}
    \hline
    \textbf{Tipo de Resíduo} & \textbf{Legislação Aplicável} \\
    \hline
    Celulares / Tablets &
      Lei 12.305/2010 (PNRS); Decreto 10.240/2020; LGPD (Lei 13.709/2018);
      Lei 9.605/1998 \\
    \hline
    Computadores / Notebooks &
      Lei 12.305/2010 (PNRS); Decreto 10.240/2020; LGPD; Lei 9.605/1998 \\
    \hline
    Baterias / Pilhas &
      Resolução CONAMA 401/2008; Lei 12.305/2010; Lei 9.605/1998 \\
    \hline
    Eletrodomésticos &
      Lei 12.305/2010 (PNRS); Decreto 10.240/2020; Lei 9.605/1998 \\
    \hline
    Outros Eletrônicos &
      Lei 12.305/2010 (PNRS); Lei 9.605/1998 \\
    \hline
  \end{tabularx}
  \fonte{Elaborado pelos autores (2025).}
\end{table}

\subsection{Modelagem do Banco de Dados}

A persistência de dados do E-Descarte adota uma estratégia híbrida, combinando
um banco relacional (PostgreSQL) para dados estruturados e de alta consistência,
e um banco de documentos (MongoDB) para dados semiestruturados, como registros
de logs de eventos e metadados de arquivos de mídia.

No PostgreSQL, as três entidades principais são:

\begin{itemize}
  \item \textbf{usuarios}: armazena credenciais (hash bcrypt), papel
        (\textit{cidadão}, \textit{gestor}, \textit{admin}) e timestamps;
  \item \textbf{pontos\_coleta}: armazena localização (latitude/longitude),
        tipo, itens aceitos e horários de funcionamento;
  \item \textbf{denuncias}: armazena protocolo único, tipo de resíduo,
        gravidade, geolocalização, status de resolução e os IDs das
        legislações relacionadas (armazenados como \texttt{TEXT[]}).
\end{itemize}

No MongoDB, duas coleções complementam a persistência:

\begin{itemize}
  \item \textbf{EventoLog}: registra todos os eventos do sistema (criação de
        denúncia, login, atualização de status) com timestamps e endereço IP,
        para auditoria e análise;
  \item \textbf{MidiaDenuncia}: armazena metadados das fotografias enviadas
        junto às denúncias, com referência ao protocolo correspondente.
\end{itemize}

\subsection{API REST}

A API do E-Descarte segue os princípios REST (\textit{Representational State
Transfer}) conforme definidos por \citeonline{fielding2000}. Os recursos são
organizados em cinco grupos de endpoints, conforme o Quadro~\ref{quad:endpoints}.

\begin{table}[H]
  \centering
  \caption{Endpoints principais da API E-Descarte}
  \label{quad:endpoints}
  \begin{tabularx}{\textwidth}{|l|l|X|l|}
    \hline
    \textbf{Método} & \textbf{Endpoint} & \textbf{Descrição} & \textbf{Auth} \\
    \hline
    POST & /api/auth/registro      & Cadastro de novo usuário & Não \\
    POST & /api/auth/login         & Autenticação e retorno de JWT & Não \\
    GET  & /api/auth/me            & Perfil do usuário autenticado & JWT \\
    \hline
    GET  & /api/pontos-coleta      & Lista pontos (filtrável por cidade/raio) & Não \\
    GET  & /api/pontos-coleta/:id  & Detalhe de um ponto & Não \\
    POST & /api/pontos-coleta      & Cria ponto de coleta & Gestor \\
    \hline
    POST & /api/denuncias          & Registra denúncia (anon. opcional) & Opcional \\
    GET  & /api/denuncias/protocolo/:p & Consulta pública por protocolo & Não \\
    PATCH& /api/denuncias/:id/status & Atualiza status & Gestor \\
    POST & /api/denuncias/:p/midia & Upload de foto & Não \\
    \hline
    GET  & /api/legislacao         & Lista toda legislação & Não \\
    GET  & /api/legislacao/tipo/:t & Legislação por tipo de resíduo & Não \\
    \hline
    POST & /api/calculadora        & Calcula impacto ambiental & Não \\
    \hline
  \end{tabularx}
  \fonte{Elaborado pelos autores (2025).}
\end{table}

\subsection{Segurança e Autenticação}

A segurança do sistema foi implementada em múltiplas camadas. A autenticação
utiliza JSON Web Tokens (JWT) \cite{jones2015}, com assinatura HMAC-SHA256 e
prazo de expiração configurável. As senhas dos usuários são armazenadas
exclusivamente sob \textit{hash} bcrypt com fator de custo 10, garantindo
resistência a ataques de dicionário e força bruta \cite{provos1999}.

Três níveis de controle de acesso foram implementados: (i) rotas públicas,
acessíveis sem autenticação (mapa de pontos, consulta de protocolo, módulo
legal); (ii) rotas com autenticação opcional, onde a identidade do usuário é
associada à denúncia quando disponível, mas o anonimato é preservado quando
solicitado; e (iii) rotas restritas a gestores e administradores (cadastro de
pontos de coleta, atualização de status de denúncias).

\subsection{Calculadora de Impacto Ambiental}

A calculadora de impacto ambiental do E-Descarte fundamenta-se em coeficientes
de impacto derivados de dados publicados pela Associação Brasileira de Limpeza
Pública e Resíduos Especiais (ABRELPE) e pela \textit{United States Environmental
Protection Agency} (EPA). Para cada categoria de dispositivo, foram definidos
quatro indicadores de impacto do descarte inadequado, conforme a
Tabela~\ref{tab:coeficientes}.

\begin{table}[H]
  \centering
  \caption{Coeficientes de impacto ambiental por tipo de dispositivo}
  \label{tab:coeficientes}
  \begin{tabular}{|l|c|c|c|c|}
    \hline
    \textbf{Dispositivo} &
    \textbf{Solo (m²)} &
    \textbf{Metais (g)} &
    \textbf{Água (L)} &
    \textbf{CO$_2$ (kg)} \\
    \hline
    Celular/Smartphone  & 0,10 & 60      & 500   & 0,5 \\
    Notebook            & 2,00 & 1.800   & 2.000 & 4,0 \\
    Bateria/Pilha       & 1,00 & 8       & 200   & 0,1 \\
    TV/Monitor          & 3,00 & 4.000   & 1.500 & 6,0 \\
    \hline
  \end{tabular}
  \fonte{Adaptado de \citeonline{abrelpe2023} e \citeonline{epa2021}.}
\end{table}

O impacto total é calculado por meio da seguinte expressão:

\begin{equation}
  I_k = \sum_{i \in D} q_i \cdot c_{i,k}
  \label{eq:impacto}
\end{equation}

\noindent
onde $I_k$ representa o impacto total no indicador $k$ (solo, metais, água ou
CO$_2$), $q_i$ é a quantidade de dispositivos do tipo $i$ a serem descartados,
e $c_{i,k}$ é o coeficiente de impacto do dispositivo $i$ no indicador $k$.

% ============================================================
\section{TECNOLOGIAS UTILIZADAS}
% ============================================================

A seleção do conjunto tecnológico do E-Descarte foi orientada pelos critérios
de maturidade do ecossistema, popularidade na comunidade de desenvolvimento,
adequação ao caso de uso e facilidade de integração entre os componentes.

\subsection{Backend}

\begin{itemize}
  \item \textbf{Node.js (v20 LTS)}: plataforma de execução JavaScript
        assíncrona e orientada a eventos, escolhida pela sua eficiência no
        tratamento de requisições I/O-intensivas e pelo compartilhamento da
        linguagem JavaScript com o \textit{frontend}, reduzindo a carga
        cognitiva da equipe \cite{tilkov2010};
  \item \textbf{Express.js (v4)}: \textit{framework} minimalista para
        construção de APIs REST em Node.js, com ecossistema de
        \textit{middlewares} maduro;
  \item \textbf{JSON Web Token (JWT)}: padrão aberto (RFC 7519) para
        transmissão segura de informações entre partes como objeto JSON
        assinado digitalmente \cite{jones2015};
  \item \textbf{bcryptjs}: implementação do algoritmo bcrypt para
        \textit{hashing} seguro de senhas \cite{provos1999};
  \item \textbf{PostgreSQL (v16)}: sistema gerenciador de banco de dados
        relacional objeto (SGBDR-O) de código aberto, escolhido pela sua
        robustez, suporte a tipos de dados avançados (como o tipo
        \texttt{TEXT[]} para arrays de strings) e conformidade com o
        padrão SQL;
  \item \textbf{MongoDB (v7)}: banco de dados orientado a documentos,
        adotado para persistência de dados semiestruturados (logs, metadados
        de mídia) por sua flexibilidade de esquema e alto desempenho em
        operações de escrita;
  \item \textbf{Mongoose (v8)}: ODM (\textit{Object Document Mapper}) para
        MongoDB em Node.js, que fornece validação de esquema e interface
        orientada a objetos;
  \item \textbf{Multer}: \textit{middleware} para tratamento de
        \textit{multipart/form-data}, utilizado no recebimento de fotografias
        das denúncias;
  \item \textbf{Nodemailer}: biblioteca para envio de e-mails transacionais
        via SMTP, utilizada nas notificações de confirmação de denúncias.
\end{itemize}

\subsection{Frontend Web}

\begin{itemize}
  \item \textbf{React (v18)}: biblioteca JavaScript declarativa para
        construção de interfaces de usuário baseadas em componentes,
        mantida pelo Meta \cite{facebook2013};
  \item \textbf{Vite (v5)}: ferramenta de \textit{build} de nova geração
        para projetos frontend, baseada em ES Modules nativo, com tempo
        de inicialização significativamente superior ao Webpack;
  \item \textbf{React Router DOM (v6)}: biblioteca de roteamento do lado
        do cliente para aplicações React de página única (SPA);
  \item \textbf{Axios}: cliente HTTP baseado em Promises para o navegador
        e Node.js, utilizado para comunicação com a API REST;
  \item \textbf{@react-google-maps/api}: \textit{wrapper} React para
        a JavaScript API do Google Maps, utilizado no mapa interativo de
        pontos de coleta.
\end{itemize}

\subsection{Aplicativo Móvel}

\begin{itemize}
  \item \textbf{React Native (v0.74)}: \textit{framework} para construção
        de aplicativos móveis nativos usando React e JavaScript, mantido
        pelo Meta. Permite o reaproveitamento de lógica de negócio e
        serviços entre o app móvel e o frontend web \cite{eisenman2015};
  \item \textbf{Expo SDK 51}: plataforma que simplifica o desenvolvimento,
        \textit{build} e distribuição de aplicativos React Native,
        fornecendo APIs para acesso a recursos nativos como câmera,
        GPS e notificações push;
  \item \textbf{expo-location}: módulo Expo para acesso à localização GPS
        do dispositivo, utilizado no georreferenciamento de denúncias;
  \item \textbf{react-native-maps}: componente de mapa nativo para React
        Native, com suporte ao Google Maps e Apple Maps;
  \item \textbf{@react-navigation}: biblioteca de navegação para React
        Native, adotando o padrão de abas inferiores (\textit{bottom tabs})
        para a navegação principal do aplicativo.
\end{itemize}

\subsection{Integração com Serviços Externos}

\begin{itemize}
  \item \textbf{Google Maps Platform}: utilizada para (i) exibição do
        mapa interativo de pontos de coleta; e (ii) geocodificação de
        endereços informados nas denúncias, convertendo endereços textuais
        em coordenadas geográficas via Geocoding API;
  \item \textbf{Twilio (ou equivalente)}: serviço de envio de SMS para
        notificação de usuários sobre o protocolo de suas denúncias.
\end{itemize}

% ============================================================
\section{RESULTADOS OBTIDOS E ESPERADOS}
% ============================================================

\subsection{Protótipo Funcional}

Ao término do ciclo de desenvolvimento descrito neste trabalho, obteve-se um
protótipo funcional completo do E-Descarte, composto por:

\begin{itemize}
  \item API REST com 15 endpoints documentados, cobrindo autenticação, CRUD
        de pontos de coleta, gerenciamento de denúncias, módulo legal e
        calculadora de impacto;
  \item Aplicação web responsiva desenvolvida em React, com cinco páginas
        funcionais: Mapa de Coleta, Registro de Denúncia, Calculadora de
        Impacto, Módulo Legal e Login/Cadastro;
  \item Aplicativo móvel desenvolvido em React Native/Expo, com quatro telas
        organizadas em navegação por abas: Mapa, Denúncia, Calculadora e
        Legislação;
  \item Banco de dados PostgreSQL inicializado com script de migração
        automatizado e dados iniciais de três pontos de coleta em Salvador (BA);
  \item Repositório público no GitHub
        (\url{https://github.com/RenoirAugusteSS/E-Descarte}) com
        estrutura de projeto organizada por camadas e documentação no
        arquivo \texttt{README.md}.
\end{itemize}

\subsection{Resultados Esperados em Implantação Real}

Embora a plataforma ainda não tenha sido implantada em ambiente de produção
para uso público em larga escala, é possível projetar os resultados esperados
com base nos dados da literatura e em analogias com sistemas similares.

\citeonline{reno2020} verificaram, em estudo conduzido com usuários de
plataformas de localização de pontos de coleta, que a disponibilização de
informações geolocalizadas aumentou em 34\% a taxa de descarte correto de
eletrônicos entre usuários da plataforma em comparação ao grupo de controle.
Extrapolando esse dado para a realidade brasileira, onde apenas 3\% dos REEE
recebem destinação formal \cite{abrelpe2023}, mesmo ganhos percentuais modestos
representariam impacto ambiental significativo dado o volume absoluto gerado.

Em relação ao Módulo Legal Integrado, a hipótese central do projeto é de que
a apresentação contextualizada da legislação ambiental -- diretamente relacionada
ao tipo de resíduo que o usuário pretende reportar -- eleva o nível de
conhecimento legal ambiental dos cidadãos usuários da plataforma, contribuindo
para o fortalecimento do controle social sobre o cumprimento da PNRS.
Essa hipótese, que requer investigação empírica sistemática para validação,
fundamenta-se na Teoria da Aprendizagem Significativa de Ausubel, segundo a
qual o aprendizado é potencializado quando o novo conhecimento é apresentado
em contexto de relevância imediata para o aprendiz \cite{ausubel2000}.

A Tabela~\ref{tab:impacto_estimado} projeta o impacto ambiental de diferentes
cenários de uso da plataforma.

\begin{table}[H]
  \centering
  \caption{Projeção de impacto ambiental por cenário de uso}
  \label{tab:impacto_estimado}
  \begin{tabular}{|l|c|c|c|}
    \hline
    \textbf{Cenário} &
    \textbf{Usuários/mês} &
    \textbf{Dispositivos descar.} &
    \textbf{Metais evitados (kg)} \\
    \hline
    Conservador   & 500    & 1.500  & 90   \\
    Moderado      & 2.000  & 6.000  & 360  \\
    Otimista      & 10.000 & 30.000 & 1.800 \\
    \hline
  \end{tabular}
  \fonte{Elaborado pelos autores com base em \citeonline{epa2021}.}
\end{table}

% ============================================================
\section{DISCUSSÃO}
% ============================================================

\subsection{Contribuições Técnicas}

Do ponto de vista da Engenharia de \textit{Software}, o E-Descarte demonstra
a viabilidade de integração entre tecnologias modernas de desenvolvimento
\textit{full-stack} (React, Node.js, PostgreSQL, MongoDB) em um sistema único
com domínio de aplicação ambiental. A adoção de persistência híbrida --
combinando banco relacional para dados transacionais e banco de documentos para
dados semiestruturados -- constitui solução elegante para os diferentes perfis
de dado gerados pelo sistema, evitando tanto a rigidez excessiva de um esquema
puramente relacional quanto a falta de garantias transacionais de uma solução
puramente NoSQL.

A implementação do Módulo Legal como componente desacoplado e reutilizável
tanto no \textit{backend} quanto nos \textit{frontends} web e móvel representa
contribuição para a engenharia de sistemas de informação com conteúdo jurídico,
demonstrando um padrão de implementação que pode ser replicado em outros
domínios de aplicação onde a correlação entre dados do usuário e normas legais
seja relevante.

\subsection{Contribuições Socioambientais}

O E-Descarte posiciona-se na interseção entre a Ciência da Computação e o
Direito Ambiental, respondendo ao chamado da EcoTech por soluções que não
apenas usem tecnologia para resolver problemas ambientais, mas que também
fortaleçam a consciência legal e cidadã dos usuários. A apresentação contextual
da legislação ambiental como funcionalidade central -- e não como simples texto
informativo em página de ``Sobre'' -- é uma escolha de design com implicações
práticas significativas: ao saber que o descarte irregular de baterias viola a
Resolução CONAMA 401/2008 e pode configurar crime ambiental pela Lei nº
9.605/1998, o cidadão ganha informação concreta para exigir o cumprimento da
lei de fabricantes, importadores e do poder público.

\subsection{Limitações}

O presente trabalho apresenta limitações que devem ser reconhecidas. Primeiro,
a ausência de avaliação empírica com usuários reais impede a validação das
hipóteses sobre o efeito do Módulo Legal na consciência legal ambiental dos
usuários. Segundo, o banco de dados de pontos de coleta iniciou-se com dados
de apenas três pontos em Salvador (BA), tornando a plataforma dependente de
alimentação de dados por gestores cadastrados para atingir cobertura adequada.
Terceiro, a ausência de integração formal com sistemas de fiscalização municipal
ou estadual limita o ciclo de resolução das denúncias registradas, que dependem
de ação voluntária dos órgãos competentes.

\subsection{Desafios Encontrados}

Os principais desafios técnicos enfrentados durante o desenvolvimento incluíram:
(i) a configuração da persistência híbrida PostgreSQL/MongoDB em ambiente de
desenvolvimento único, exigindo cuidados especiais na gestão de conexões e no
tratamento de erros de cada banco de forma independente; (ii) a implementação
do upload de arquivos com metadados persistidos em banco de dados diferente do
arquivo físico, resolvida pelo armazenamento dos metadados no MongoDB e do
arquivo no sistema de arquivos local (com perspectiva de migração para
armazenamento em nuvem em versão futura); e (iii) a garantia do anonimato nas
denúncias, que exigiu cuidado especial no \textit{middleware} de autenticação
opcional para evitar que dados de identidade fossem inadvertidamente associados
a denúncias anônimas.

% ============================================================
\section{CONCLUSÃO}
% ============================================================

Este artigo apresentou o E-Descarte, uma plataforma digital desenvolvida para
a gestão e denúncia de resíduos eletroeletrônicos no contexto brasileiro,
com ênfase no Módulo Legal Integrado como diferencial sociotécnico. O sistema
atende ao cenário crítico de geração massiva de REEE no Brasil -- 2,1 milhões
de toneladas anuais, com menos de 3\% de destinação formal adequada -- por meio
de uma solução tecnológica que combina geolocalização, controle social e
educação legal ambiental em uma única plataforma acessível via web e dispositivo
móvel.

Os objetivos propostos foram integralmente alcançados: desenvolveu-se uma API
REST completa com 15 endpoints, uma aplicação web React com cinco páginas
funcionais, um aplicativo React Native/Expo com quatro telas e um Módulo Legal
Integrado que mapeia dinamicamente os tipos de resíduo eletroeletrônico à
legislação ambiental brasileira aplicável (Lei 12.305/2010, Decreto 10.240/2020,
Resolução CONAMA 401/2008, Lei 9.605/1998 e Lei 13.709/2018).

Do ponto de vista ambiental, o E-Descarte contribui para a redução do hiato
entre a robustez do ordenamento jurídico ambiental brasileiro e a efetividade
de sua aplicação, instrumentalizando o controle social por meio da informação
geolocalizada e contextualizada. Do ponto de vista tecnológico, demonstra a
viabilidade e adequação das tecnologias Node.js, React, React Native, PostgreSQL
e MongoDB para o desenvolvimento de sistemas de informação ambiental modernos,
escaláveis e acessíveis.

O E-Descarte representa, ainda, um protótipo funcional para integração no
Observatório Digital Ambiental da Turma, compondo -- junto a outros projetos
da disciplina EcoTech -- um panorama tecnológico das soluções que a Computação
pode oferecer ao enfrentamento dos desafios ambientais contemporâneos.

% ============================================================
\section{TRABALHOS FUTUROS}
% ============================================================

A partir da versão atual do protótipo, identificam-se as seguintes direções
prioritárias para trabalhos futuros:

\begin{enumerate}
  \item \textbf{Avaliação de usabilidade e impacto}: conduzir estudo empírico
        com usuários reais para avaliar a usabilidade da plataforma (escala
        SUS -- \textit{System Usability Scale}) e o efeito do Módulo Legal
        Integrado no nível de conhecimento legal ambiental dos usuários (pré e
        pós-teste);

  \item \textbf{Integração com órgãos de fiscalização}: desenvolver interfaces
        de integração com os sistemas de ouvidoria e fiscalização ambiental de
        municípios e estados, criando um ciclo completo de resolução para as
        denúncias registradas na plataforma;

  \item \textbf{Crowdsourcing de pontos de coleta}: implementar funcionalidade
        que permita ao próprio cidadão sugerir novos pontos de coleta, com
        moderação por gestores cadastrados, acelerando a expansão da cobertura
        geográfica da plataforma;

  \item \textbf{Gamificação e programa de fidelidade verde}: incorporar
        mecanismos de gamificação (pontos, medalhas, rankings) para incentivar
        o uso recorrente da plataforma e premiar cidadãos com histórico
        consistente de descarte correto;

  \item \textbf{Dashboard analítico para gestores públicos}: desenvolver painel
        de controle com visualizações de dados sobre concentração geográfica de
        denúncias, tendências temporais e indicadores de resolução, auxiliando
        na tomada de decisão dos órgãos ambientais municipais e estaduais;

  \item \textbf{Análise preditiva com Machine Learning}: aplicar algoritmos de
        agrupamento (\textit{clustering}) sobre os dados históricos de denúncias
        para identificar áreas de risco de descarte irregular, permitindo ação
        preventiva dos órgãos de fiscalização;

  \item \textbf{Armazenamento em nuvem e escalabilidade}: migrar o armazenamento
        de arquivos de mídia para serviços de armazenamento em nuvem (AWS S3,
        Google Cloud Storage) e containerizar a aplicação com Docker e
        Kubernetes para suportar escalonamento em produção;

  \item \textbf{Atualização dinâmica do módulo legal}: integrar o módulo legal
        a feeds jurídicos automatizados (como o Diário Oficial da União) para
        que alterações legislativas sejam refletidas na plataforma sem
        necessidade de intervenção manual.
\end{enumerate}

% ============================================================
%  REFERÊNCIAS BIBLIOGRÁFICAS
% ============================================================
\clearpage
\begin{center}
  {\normalsize\bfseries REFERÊNCIAS}
\end{center}

\begin{thebibliography}{99}

\bibitem{abrelpe2023}
ASSOCIAÇÃO BRASILEIRA DE LIMPEZA PÚBLICA E RESÍDUOS ESPECIAIS -- ABRELPE.
\textbf{Panorama dos Resíduos Sólidos no Brasil 2022}. São Paulo: ABRELPE,
2023. Disponível em: \url{https://abrelpe.org.br/panorama/}. Acesso em:
10 jun. 2025.

\bibitem{ausubel2000}
AUSUBEL, D. P.
\textbf{The acquisition and retention of knowledge: a cognitive view}.
Dordrecht: Kluwer Academic Publishers, 2000.

\bibitem{baldé2017}
BALDÉ, C. P. et al.
\textbf{The Global E-waste Monitor 2017}: quantities, flows and resources.
Bonn: United Nations University, International Telecommunication Union,
International Solid Waste Association, 2017.

\bibitem{brasil1998crimes}
BRASIL. Lei nº 9.605, de 12 de fevereiro de 1998. Dispõe sobre as sanções
penais e administrativas derivadas de condutas e atividades lesivas ao meio
ambiente, e dá outras providências.
\textbf{Diário Oficial da União}, Brasília, DF, 13 fev. 1998. Disponível em:
\url{https://www.planalto.gov.br/ccivil_03/leis/l9605.htm}. Acesso em:
10 jun. 2025.

\bibitem{brasil2010pnrs}
BRASIL. Lei nº 12.305, de 02 de agosto de 2010. Institui a Política Nacional
de Resíduos Sólidos; altera a Lei nº 9.605, de 12 de fevereiro de 1998; e dá
outras providências.
\textbf{Diário Oficial da União}, Brasília, DF, 03 ago. 2010. Disponível em:
\url{https://www.planalto.gov.br/ccivil_03/_ato2007-2010/2010/lei/l12305.htm}.
Acesso em: 10 jun. 2025.

\bibitem{brasil2018lgpd}
BRASIL. Lei nº 13.709, de 14 de agosto de 2018. Lei Geral de Proteção de
Dados Pessoais (LGPD).
\textbf{Diário Oficial da União}, Brasília, DF, 15 ago. 2018. Disponível em:
\url{https://www.planalto.gov.br/ccivil_03/_ato2015-2018/2018/lei/l13709.htm}.
Acesso em: 10 jun. 2025.

\bibitem{brasil2020decreto}
BRASIL. Decreto nº 10.240, de 12 de fevereiro de 2020. Regulamenta o inciso VI
do caput do art. 33 da Lei nº 12.305, de 2 de agosto de 2010, e o § 1º do art.
56 do Decreto nº 9.177, de 23 de outubro de 2017.
\textbf{Diário Oficial da União}, Brasília, DF, 12 fev. 2020. Disponível em:
\url{https://www.planalto.gov.br/ccivil_03/_ato2019-2022/2020/decreto/D10240.htm}.
Acesso em: 10 jun. 2025.

\bibitem{conama401}
CONSELHO NACIONAL DO MEIO AMBIENTE -- CONAMA. Resolução nº 401, de 04 de
novembro de 2008. Estabelece os limites máximos de chumbo, cádmio e mercúrio
para pilhas e baterias comercializadas no território nacional.
\textbf{Diário Oficial da União}, Brasília, DF, 05 nov. 2008.

\bibitem{earth911}
EARTH911. \textbf{iRecycle: America's recycling app}. [S.l.]: Earth911, 2023.
Disponível em: \url{https://earth911.com}. Acesso em: 10 jun. 2025.

\bibitem{eisenman2015}
EISENMAN, B.
\textbf{Learning React Native: building native mobile apps with JavaScript}.
Sebastopol: O'Reilly Media, 2015.

\bibitem{epa2021}
UNITED STATES ENVIRONMENTAL PROTECTION AGENCY -- EPA.
\textbf{Advancing Sustainable Materials Management: 2019 fact sheet}.
Washington, D.C.: EPA, 2021. Disponível em:
\url{https://www.epa.gov/facts-and-figures-about-materials-waste-and-recycling}.
Acesso em: 10 jun. 2025.

\bibitem{europarlamento2012}
PARLAMENTO EUROPEU. Diretiva 2012/19/UE do Parlamento Europeu e do Conselho,
de 4 de julho de 2012, relativa aos resíduos de equipamentos elétricos e
eletrônicos (REEE).
\textbf{Jornal Oficial da União Europeia}, Bruxelas, 24 jul. 2012.

\bibitem{facebook2013}
FACEBOOK. \textbf{React: a JavaScript library for building user interfaces}.
[S.l.]: Meta Platforms, 2013. Disponível em: \url{https://react.dev}.
Acesso em: 10 jun. 2025.

\bibitem{fielding2000}
FIELDING, R. T.
\textbf{Architectural styles and the design of network-based software
architectures}. 2000. Dissertação (Doutorado) -- University of California,
Irvine, 2000.

\bibitem{forti2020}
FORTI, V. et al.
\textbf{The Global E-waste Monitor 2020}: quantities, flows and the circular
economy potential. Bonn: United Nations University; United Nations Institute
for Training and Research; International Telecommunication Union; International
Solid Waste Association, 2020.

\bibitem{fowler2014}
FOWLER, M.; LEWIS, J. \textbf{Microservices}: a definition of this new
architectural term. [S.l.]: martinfowler.com, 2014. Disponível em:
\url{https://martinfowler.com/articles/microservices.html}.
Acesso em: 10 jun. 2025.

\bibitem{gil2017}
GIL, A. C.
\textbf{Como elaborar projetos de pesquisa}. 6. ed. São Paulo: Atlas, 2017.

\bibitem{ibama2022}
INSTITUTO BRASILEIRO DO MEIO AMBIENTE E DOS RECURSOS NATURAIS RENOVÁVEIS --
IBAMA.
\textbf{Relatório de qualidade do meio ambiente}. Brasília: IBAMA, 2022.
Disponível em: \url{https://www.ibama.gov.br}. Acesso em: 10 jun. 2025.

\bibitem{jones2015}
JONES, M. B. et al.
\textbf{JSON Web Token (JWT): RFC 7519}. [S.l.]: Internet Engineering Task
Force, 2015. Disponível em: \url{https://www.rfc-editor.org/rfc/rfc7519}.
Acesso em: 10 jun. 2025.

\bibitem{moreira2019}
MOREIRA, A. B.; FIGUEIREDO, M. L.
Gestão de resíduos eletroeletrônicos no Brasil: análise do marco regulatório
e desafios de implementação da logística reversa.
\textbf{Revista de Gestão Ambiental e Sustentabilidade}, São Paulo, v. 8,
n. 2, p. 310--328, 2019.

\bibitem{provos1999}
PROVOS, N.; MAZIÈRES, D.
A future-adaptable password scheme.
In: \textbf{Proceedings of the Annual Technical Conference on USENIX}, 1999,
Monterey. \textbf{Anais}[\ldots] Berkeley: USENIX Association, 1999.
p. 81--91.

\bibitem{reno2020}
RENO, M. A. et al.
Impact of geolocation-based recycling apps on electronic waste disposal
behavior: a randomized controlled study.
\textbf{Journal of Cleaner Production}, Amsterdam, v. 276, p. 124--135,
dez. 2020.

\bibitem{ribeiro2021}
RIBEIRO, F. M.; SANTOS, D. C.; OLIVEIRA, T. P.
Aplicativos móveis para gestão de resíduos sólidos no Brasil: um mapeamento
sistemático da literatura.
\textbf{Revista Brasileira de Ciências Ambientais}, Rio de Janeiro, v. 56,
n. 1, p. 1--16, 2021.

\bibitem{singh2020}
SINGH, N. et al.
Electronic waste (e-waste): a global threat to the environment and human health.
\textbf{Waste Management \& Research}, London, v. 38, n. 1, p. 3--24, 2020.

\bibitem{sommerville2019}
SOMMERVILLE, I.
\textbf{Engenharia de software}. 10. ed. São Paulo: Pearson Education, 2019.

\bibitem{tilkov2010}
TILKOV, S.; VINOSKI, S.
Node.js: using JavaScript to build high-performance network programs.
\textbf{IEEE Internet Computing}, v. 14, n. 6, p. 80--83, nov./dez. 2010.

\bibitem{trashout}
TRASHOUT. \textbf{TrashOut: the worldwide trash reporting service}.
[S.l.]: TrashOut s.r.o., 2023. Disponível em: \url{https://www.trashout.ngo}.
Acesso em: 10 jun. 2025.

\bibitem{watson2010}
WATSON, R. T. et al.
\textbf{Information systems and environmentally sustainable development}:
energy informatics and new directions for the IS community.
\textbf{MIS Quarterly}, Minneapolis, v. 34, n. 1, p. 23--38, mar. 2010.

\end{thebibliography}

\end{document}
